% GENERATED FILE. Edit canonical lessons, then run npm run generate:books.
% canonical-source-hash: fnv1a64:5128ccbceee3fcaa
% canonical-lessons: MR-C472-hatti, MR-C472-admutha, MR-C472-lahri, MR-C472-utavla, MR-C472-adhir

\chapter{Stubborn and Hasty}
\label{ch:mr-stubborn-and-hasty}
\hlchaptermodality{472}

\begin{chapteropening}
\textbf{By the end of this chapter:} \emph{I can describe a stubborn, moody, hasty or impatient person.}

Complete the last lesson of chapter 472: I can describe a stubborn, moody, hasty or impatient person.
\end{chapteropening}

\section[\texorpdfstring{\mr{हट्टी} (haṭṭī)}{हट्टी (haṭṭī)}]{\mr{हट्टी} (haṭṭī) — stubborn}

\label{lesson:MR-C472-hatti}



\begin{quote}
Before the new one: say the Marathi for understanding, then the Marathi for common sense.
\end{quote}

\subsection*{You'll want to know: \mr{हट्टी}}
\textbf{\mr{हट्टी}} — \emph{haṭṭī} — \textquotedblleft{}stubborn\textquotedblright{}.

\begin{quote}
\mr{हट्टी} \mr{बाळ} \mr{खेळण्यासाठी} \mr{रडत} \mr{राहिलं}.
\end{quote}

\begin{quote}
\emph{haṭṭī bāḷ kheḷṇyāsāṭhī raḍat rāhilaṁ.}
\end{quote}

\begin{quote}
The stubborn baby kept crying for the toy.
\end{quote}

From \mr{हट्ट}, insistence: \mr{हट्ट} \mr{करणे}, to insist like a child. Unchanging.

\begin{cousinweb}
\textbf{\mr{हट्टी}} (\emph{haṭṭī}) is a Marathi adjective built on \mr{हट्ट}, \textquotedblleft{}insistence\textquotedblright{}, which goes back to Sanskrit \emph{haṭha}, \textquotedblleft{}force; obstinacy\textquotedblright{} (Monier-Williams).
\end{cousinweb}

\begin{grammarlens}[title={What you've built}]
One more word for describing what you lived and felt.
\end{grammarlens}

\subsection*{Your turn}
\begin{itemize}
\raggedright
  \item \emph{Say it:} \emph{haṭṭī}
  \item \emph{Say it:} \emph{haṭṭī}, once more
  \item \emph{Say it:} \emph{haṭṭī bāḷ kheḷṇyāsāṭhī raḍat rāhilaṁ.} — The stubborn baby kept crying for the toy.
  \item \emph{Recall:} say the Marathi for understanding, then the Marathi for common sense, then say \emph{haṭṭī} again
\end{itemize}

\begin{tcolorbox}[breakable,colback=teal!4,colframe=teal!35!black,title={Before you move on}]
What does \mr{हट्टी} mean? (\textquotedblleft{}Stubborn\textquotedblright{}.) Say it once more.
\end{tcolorbox}
\section[\texorpdfstring{\mr{आडमुठा} (āḍmuṭhā)}{आडमुठा (āḍmuṭhā)}]{\mr{आडमुठा} (āḍmuṭhā) — obstinate; pig-headed}

\label{lesson:MR-C472-admutha}



\begin{quote}
Before the new one: say the Marathi for common sense, then the Marathi for stubborn.
\end{quote}

\subsection*{You'll want to know: \mr{आडमुठा}}
\textbf{\mr{आडमुठा}} — \emph{āḍmuṭhā} — \textquotedblleft{}obstinate; pig-headed\textquotedblright{}.

\begin{quote}
\mr{तो} \mr{इतका} \mr{आडमुठा} \mr{आहे} \mr{की} \mr{कोणाचंच} \mr{ऐकत} \mr{नाही}.
\end{quote}

\begin{quote}
\emph{to itkā āḍmuṭhā āhe kī koṇācañc aikat nāhī.}
\end{quote}

\begin{quote}
He is so pig-headed that he listens to no one.
\end{quote}

Agrees; harsher than \mr{हट्टी}.

\begin{cousinweb}
\textbf{\mr{आडमुठा}} (\emph{āḍmuṭhā}) is a Marathi adjective whose origin is not settled.
\end{cousinweb}

\begin{grammarlens}[title={What you've built}]
One more word for describing what you lived and felt.
\end{grammarlens}

\subsection*{Your turn}
\begin{itemize}
\raggedright
  \item \emph{Say it:} \emph{āḍmuṭhā}
  \item \emph{Say it:} \emph{āḍmuṭhā}, once more
  \item \emph{Say it:} \emph{to itkā āḍmuṭhā āhe kī koṇācañc aikat nāhī.} — He is so pig-headed that he listens to no one.
  \item \emph{Recall:} say the Marathi for common sense, then the Marathi for stubborn, then say \emph{āḍmuṭhā} again
\end{itemize}

\begin{tcolorbox}[breakable,colback=teal!4,colframe=teal!35!black,title={Before you move on}]
What does \mr{आडमुठा} mean? (\textquotedblleft{}Obstinate; pig-headed\textquotedblright{}.) Say it once more.
\end{tcolorbox}
\section[\texorpdfstring{\mr{लहरी} (lahrī)}{लहरी (lahrī)}]{\mr{लहरी} (lahrī) — moody; whimsical}

\label{lesson:MR-C472-lahri}



\begin{quote}
Before the new one: say the Marathi for stubborn, then the Marathi for obstinate.
\end{quote}

\subsection*{You'll want to know: \mr{लहरी}}
\textbf{\mr{लहरी}} — \emph{lahrī} — \textquotedblleft{}moody; whimsical\textquotedblright{}.

\begin{quote}
\mr{आमचे} \mr{मामा} \mr{लहरी} \mr{आहेत}, \mr{कधी} \mr{हसतात} \mr{तर} \mr{कधी} \mr{रागावतात}.
\end{quote}

\begin{quote}
\emph{āmce māmā lahrī āhet, kadhī hastāt tar kadhī rāgāvtāt.}
\end{quote}

\begin{quote}
Our uncle is moody: sometimes he laughs, sometimes he gets angry.
\end{quote}

From \mr{लहर}, a wave or a whim, + -\mr{ई}. Unchanging.

\begin{cousinweb}
\textbf{\mr{लहरी}} (\emph{lahrī}) is a Marathi adjective built on \mr{लहर}, \textquotedblleft{}wave, caprice\textquotedblright{}, which CDIAL lists under Sanskrit \emph{laharī}, \textquotedblleft{}large wave\textquotedblright{} (CDIAL 10999).
\end{cousinweb}

\begin{grammarlens}[title={What you've built}]
One more word for describing what you lived and felt.
\end{grammarlens}

\subsection*{Your turn}
\begin{itemize}
\raggedright
  \item \emph{Say it:} \emph{lahrī}
  \item \emph{Say it:} \emph{lahrī}, once more
  \item \emph{Say it:} \emph{āmce māmā lahrī āhet, kadhī hastāt tar kadhī rāgāvtāt.} — Our uncle is moody: sometimes he laughs, sometimes he gets angry.
  \item \emph{Recall:} say the Marathi for stubborn, then the Marathi for obstinate, then say \emph{lahrī} again
\end{itemize}

\begin{tcolorbox}[breakable,colback=teal!4,colframe=teal!35!black,title={Before you move on}]
What does \mr{लहरी} mean? (\textquotedblleft{}Moody; whimsical\textquotedblright{}.) Say it once more.
\end{tcolorbox}
\section[\texorpdfstring{\mr{उतावळा} (utāvḷā)}{उतावळा (utāvḷā)}]{\mr{उतावळा} (utāvḷā) — hasty; rash}

\label{lesson:MR-C472-utavla}



\begin{quote}
Before the new one: say the Marathi for obstinate, then the Marathi for moody.
\end{quote}

\subsection*{You'll want to know: \mr{उतावळा}}
\textbf{\mr{उतावळा}} — \emph{utāvḷā} — \textquotedblleft{}hasty; rash\textquotedblright{}.

\begin{quote}
\mr{उतावळा} \mr{माणूस} \mr{घाईत} \mr{चुका} \mr{करतो}.
\end{quote}

\begin{quote}
\emph{utāvḷā māṇūs ghāīt cukā karto.}
\end{quote}

\begin{quote}
A hasty person makes mistakes in a rush.
\end{quote}

Agrees: \mr{उतावळा}, \mr{उतावळी}, \mr{उतावळं}.

\begin{cousinweb}
\textbf{\mr{उतावळा}} (\emph{utāvḷā}) is inherited: CDIAL lists Marathi \emph{utāvḷā}, \textquotedblleft{}impetuous\textquotedblright{}, under the reconstructed \emph{*uttāpala} (CDIAL 1788).
\end{cousinweb}

\begin{grammarlens}[title={What you've built}]
One more word for describing what you lived and felt.
\end{grammarlens}

\subsection*{Your turn}
\begin{itemize}
\raggedright
  \item \emph{Say it:} \emph{utāvḷā}
  \item \emph{Say it:} \emph{utāvḷā}, once more
  \item \emph{Say it:} \emph{utāvḷā māṇūs ghāīt cukā karto.} — A hasty person makes mistakes in a rush.
  \item \emph{Recall:} say the Marathi for obstinate, then the Marathi for moody, then say \emph{utāvḷā} again
\end{itemize}

\begin{tcolorbox}[breakable,colback=teal!4,colframe=teal!35!black,title={Before you move on}]
What does \mr{उतावळा} mean? (\textquotedblleft{}Hasty; rash\textquotedblright{}.) Say it once more.
\end{tcolorbox}
\section[\texorpdfstring{\mr{अधीर} (adhīr)}{अधीर (adhīr)}]{\mr{अधीर} (adhīr) — impatient}

\label{lesson:MR-C472-adhir}



\begin{quote}
Before the new one: say the Marathi for moody, then the Marathi for hasty.
\end{quote}

\subsection*{You'll want to know: \mr{अधीर}}
\textbf{\mr{अधीर}} — \emph{adhīr} — \textquotedblleft{}impatient\textquotedblright{}.

\begin{quote}
\mr{निकाल} \mr{ऐकायला} \mr{सगळे} \mr{अधीर} \mr{झाले} \mr{होते}.
\end{quote}

\begin{quote}
\emph{nikāl aikāylā sagḷe adhīr jhāle hote.}
\end{quote}

\begin{quote}
Everyone was impatient to hear the result.
\end{quote}

\mr{अ}-, not, + \mr{धीर}, patience: \mr{अधीर} \mr{होणे}, to grow impatient.

\begin{cousinweb}
\textbf{\mr{अधीर}} (\emph{adhīr}) is a tatsama, taken from Sanskrit \emph{adhīra} (Monier-Williams), the negative \emph{a-} with \emph{dhīra}, \textquotedblleft{}steady\textquotedblright{}.
\end{cousinweb}

\begin{grammarlens}[title={What you've built}]
One more word for describing what you lived and felt.
\end{grammarlens}

\subsection*{Your turn}
\begin{itemize}
\raggedright
  \item \emph{Say it:} \emph{adhīr}
  \item \emph{Say it:} \emph{adhīr}, once more
  \item \emph{Say it:} \emph{nikāl aikāylā sagḷe adhīr jhāle hote.} — Everyone was impatient to hear the result.
  \item \emph{Recall:} say the Marathi for moody, then the Marathi for hasty, then say \emph{adhīr} again
\end{itemize}

\begin{tcolorbox}[breakable,colback=teal!4,colframe=teal!35!black,title={Before you move on}]
What does \mr{अधीर} mean? (\textquotedblleft{}Impatient\textquotedblright{}.) Say it once more.
\end{tcolorbox}
